\section{联想记忆专题练习：外积、干扰与容量}

\begin{enumerate}[label=\textbf{题 \arabic*.}]
\item $\vect{a}=(1,0)$，$\vect{b}=(0.8,0.6)$，两者点积是多少？相似度高还是低？
\item 如果只存一条记忆 $W=\vect{v}\vect{k}^\top$ 且 $\vect{k}$ 为单位向量，那么查询 $\vect{k}$ 会得到什么？
\item 为什么 $W=\sum_i\vect{v}_i\vect{k}_i^\top$ 在查询某条记忆时可能产生误差？
\item 比较 $0.5$、$0.5^2$、$0.5^3$：越来越高次的激活为什么有助于压低弱匹配？
\item Softmax 对两个分数 $[1,0.6]$ 做指数归一化。增大逆温度 $\beta=1/\tau$（$\tau$ 为温度）后，哪条记忆的权重上升？
\item 若 $K=\begin{bmatrix}1&2\\3&4\end{bmatrix}$、$\vect{x}=\begin{bmatrix}1\\-1\end{bmatrix}$，计算 $\relu(K\vect{x})$。
\item 若 $V^\top=\begin{bmatrix}2&0\\1&3\end{bmatrix}$、$\vect{a}=\begin{bmatrix}0\\2\end{bmatrix}$，计算 $V^\top\vect{a}$，并用「按列加权求和」解释结果。
\item 已知 $\vect a=(6,0)$、$\vect b=(5,12)$，求点积、两个长度与 $\cos\theta$，并验证 $\vect a\cdot\vect b=\|\vect a\|\|\vect b\|\cos\theta$。
\item 若 $\vect v_1=(2,1)^\top$、$\vect v_2=(0,3)^\top$、$\vect a=(0.8,0.2)^\top$，写出按行保存的 $V$ 与按列读取的 $V^\top$，计算 $V^\top\vect a$。
\item 用一句话说明为什么简化模型可以取 $\vect q=\vect x$，而真实 Transformer 中两者往往不同。
\item 若 $Q:2\times3$、$K:4\times3$、$V:4\times5$，求 $QK^\top$ 与 $\softmax(QK^\top)V$ 的尺寸；Softmax 沿哪一个方向归一化？
\end{enumerate}

\noindent\emph{参考解答统一放在\hyperref[app:answers]{附录答案}，请先独立完成。}



\subsection*{联想记忆的关键结论}
\begin{conclusion}
\noindent\emph{一张脑中的路线图}\par

\textbf{点积}：度量 Query 与 Key 的匹配程度。\\
\textbf{外积}：把 Key--Value 关联写入一个矩阵。\\
\textbf{线性召回}：把匹配项和不匹配项一起混合。\\
\textbf{MLP 非线性}：抑制弱匹配，从而缓解干扰。\\
\textbf{Softmax Attention}：把检索权重集中在更相关的记忆上。\\[2mm]
\textbf{最核心的一个等式：}
\[
\boxed{W=\sum_i\vect{v}_i\vect{k}_i^\top
\quad\Longrightarrow\quad W\vect{x}=\sum_i\vect{v}_i(\vect{k}_i^\top\vect{x}).}
\]
这说明：矩阵乘法也可以被理解为按相似度检索一组关联记忆。
\end{conclusion}
